\SourceLeaf{106}%
tung dieser Blätter nicht angemessen sein, diesen Untersuchungen hier Schritt vor Schritt zu folgen, auch unnöthig, da die Abhandlung selbst bereits unter der Presse ist. Wir begnügen uns daher, nur die Gegenstände von einigen dieser Untersuchungen, die sich leichter isolirt herausheben lassen, hier anzuführen.

Die Werthe der unbekannten Grössen, welche der Methode der kleinsten Quadrate gemäss sind. und die man die \textit{sichersten Werthe} nennen kann, werden vermittelst einer bestimmten Elimination gefunden, und die diesen Bestimmungen beizulegenden Gewichte vermittelst einer unbestimmten Elimination, wie dies schon aus der \textit{Theoria motus Corporum Coelestium} bekannt ist: auf eine neue Art wird hier \textit{a priori} bewiesen, dass unter den obwaltenden Voraussetzungen diese Elimination allemal möglich ist. Zugleich wird eine merkwürdige Symmetrie unter den bei der unbestimmten Elimination hervorgehenden Coëfficienten nachgewiesen.

So leicht und klar sich diese Eliminationsgeschäfte im Allgemeinen übersehen lassen, so ist doch nicht zu läugnen, dass die wirkliche numerische Ausführung, bei einer beträchtlichen Anzahl von unbekannten Grössen, beschwerlich wird. Was die \textit{bestimmte} Elimination, die zur Ausmittelung der sichersten Werthe für die unbekannten Grössen zureicht, betrifft, so hat der Verfasser ein Verfahren, wodurch die wirkliche Rechnung, so viel es nur die Natur der Sache verträgt, abgekürzt wird, bereits in der \textit{Theoria Motus Corporum Coelestium} angedeutet, und in einer im ersten Bande der \textit{Commentt. Rec. Soc. R. Gott.} befindlichen Abhandlung, \textit{Disquisitio de elementis ellipticis Palladis}, ausführlich entwickelt. Dieses Verfahren gewährt zugleich den Vortheil, dass das Gewicht der Bestimmung der einen unbekannten Grösse, welche man bei dem Geschäft als die letzte betrachtet hat, sich von selbst mit ergibt. Da nun die Ordnung unter den unbekannten Grössen gänzlich willkürlich ist, und man also welche man will, als die letzte behandeln kann, so ist dies Verfahren in allen Fällen zureichend, wo nur für Eine der unbekannten Grössen das Gewicht mit verlangt wird, und die beschwerliche unbestimmte Elimination wird dann umgangen.

Die seitdem bei den rechnenden Astronomen so allgemein gewordene Gewohnheit, die Methode der kleinsten Quadrate auf schwierige astronomische Rechnungen anzuwenden, wie auf die vollständige Bestimmung von Cometenbahnen, wobei die Anzahl der unbekannten Grössen bis auf sechs steigt, hat indess das Bedürfniss, das Gewicht der sichersten Werthe \textit{aller} unbekannten Grössen auf
\SourceLeaf{107}
eine bequemere Art als durch die unbestimmte Elimination, zu finden, fühlbar gemacht, und da die Bemühungen einiger Geometer\footnote{z. B. \textsc{Plana}'s. Siehe Zeitschrift für Astronomie und verwandte Wissenschaften Band 6, S. 258.} keinen Erfolg gehabt hatten, so hat man sich nur so geholfen, dass man den oben erwähnten Algorithmus so viele male mit veränderter Ordnung der unbekannten Grössen durchführte, als unbekannte Grössen waren, indem man jeder einmal den letzten Platz anwies. Es scheint uns jedoch, dass durch dieses kunstlose Verfahren in Vergleichung mit der unbestimmten Elimination in Rücksicht auf Kürze der Rechnung nichts gewonnen wird. Der Verfasser hat daher diesen wichtigen Gegenstand einer besondern Untersuchung unterworfen, und einen neuen Algorithmus zur Bestimmung der Gewichte der Werthe \textit{sämmtlicher} unbekannten Grössen mitgetheilt. der alle Geschmeidigkeit und Kürze zu haben scheint, welcher die Sache ihrer Natur nach fähig ist.

Der sicherste Werth einer Grösse, welche eine gegebene Function der unbekannten Grössen der Aufgabe ist, wird gefunden, indem man für letztere ihre durch die Methode der kleinsten Quadrate erhaltenen sichersten Werthe substituirt. Allein eine bisher noch nicht behandelte Aufgabe ist es, wie das jener Bestimmung beizulegende Gewicht zu finden sei. Die hier gegebene Auflösung dieser Aufgabe verdient um so mehr von den rechnenden Astronomen beherzigt zu werden, da sich findet, dass mehrere derselben dabei früher auf eine nicht richtige Art zu Werke gegangen sind.

Die Summe der Quadrate der Unterschiede zwischen den unmittelbar beobachteten Grössen, und denjenigen Werthen, welchen ihre Ausdrücke, als Functionen der unbekannten Grössen, durch Substitution der sichersten Werthe für letztere erhalten (welche Quadrate, im Fall die Beobachtungen ungleiche Zuverlässigkeit haben, vor der Addition erst noch durch die respectiven Gewichte multiplicirt werden müssen) bildet bekanntlich ein absolutes Minimum. Sobald man daher einer der unbekannten Grössen einen Werth beilegt, der von dem sichersten verschieden ist, wird ein ähnliches Aggregat, wie man auch die übrigen unbekannten Grössen bestimmen mag, allezeit grösser ausfallen, als das erwähnte Minimum. Allein die übrigen unbekannten Grössen werden sich nur auf Eine Art so bestimmen lassen, dass die Vergrösserung des Aggregats so klein wie möglich, oder dass das Aggregat selbst ein relatives Minimum werde. Diese von dem Ver\SourceLeaf{108}fasser hier ausgeführte Untersuchung führt. zu einigen interessanten Wahrheiten, die über die ganze Lehre noch ein vielseitigeres Licht verbreiten.

Es fügt sich zuweilen, dass man erst, nachdem man schon eine ausgedehnte Rechnung über eine Reihe von Beobachtungen in allen Theilen durchgeführt hat, Kenntniss von einer neuen Beobachtung erhält, die man gern noch mit zugezogen hätte. Es kann in vielen Fällen erwünscht sein, wenn man nicht nöthig hat, deshalb die ganze Eliminationsarbeit von vorne wieder anzufangen, sondern im Stande ist, die durch das Hinzukommen der neuen Beobachtung entstehende Modification in den sichersten Werthen und deren Gewichten zu finden. Der Verfasser hat daher diese Aufgabe hier besonders abgehandelt, eben so wie die verwandte, wo man einer schon angewandten Beobachtung hintennach ein anderes Gewicht, als ihr beigelegt war, zu ertheilen sich veranlasst sieht, und, ohne die Rechnung von vorne zu wiederholen, die Veränderungen der Endresultate zu erhalten wünscht.

Wie der \textit{wahrscheinliche} Fehler einer Beobachtungsgattung (als bisher üblicher Maassstab ihrer Unsicherheit) aus einer hinlänglichen Anzahl wirklicher Beobachtungsfehler näherungsweise zu finden sei, hatte der Verfasser in einer besondern Abhandlung in der Zeitschrift für Astronomie und verwandte Wissenschaften [1816. März u. April] gezeigt: dieses Verfahren, so wie der Gebrauch des wahrscheinlichen Fehlers überhaupt, ist aber von der hypothetischen Form der Grösse der Wahrscheinlichkeit der einzelnen Fehler abhängig, und musste es sein. Im ersten Theile der gegenwärtigen Abhandlung ist nun zwar gezeigt, wie aus denselben Datis der mittlere Fehler der Beobachtungen (als zweckmässiger Maassstab ihrer Ungenauigkeit) näherungsweise gefunden wird. Allein immer bleibt hiebei die Bedenklichkeit übrig, dass man nach aller Schärfe selten oder fast nie im Besitz der Kenntniss der wahren Grösse von einer Anzahl wirklicher Beobachtungsfehler sein kann. Bei der Ausübung hat man dafür bisher immer die Unterschiede zwischen dem, was die Beobachtungen ergeben haben und dem Resultaten der Rechnung nach den durch die Methode der kleinsten Quadrate gefundenen sichersten Werthen der unbekannten Grössen, wovon die Beobachtungen abhangen, zum Grunde gelegt. Allein da man nicht berechtigt ist, die sichersten Werthe für die wahren Werthe selbst zu halten, so überzeugt man sich leicht, dass man durch dieses Verfahren allemal den wahrscheinlichen und mittlern Fehler \textit{zu klein} finden muss, und daher den Beobachtungen und den daraus gezoge%
